% SPDX-License-Identifier: CC-BY-4.0
\section{The \chronosynd{} system}
\label{sec:system}

\subsection{Pipeline}

\chronosynd{} is a one-directional pipeline:
\begin{equation*}
  \text{kernel BPF}
  \,\longrightarrow\,
  \text{ring buffer}
  \,\longrightarrow\,
  \text{collector}
  \,\longrightarrow\,
  \text{features}
  \,\longrightarrow\,
  \text{scorer}
  \,\longrightarrow\,
  \text{alerts}
\end{equation*}
Kernel-side BPF programs emit raw syscall events into a userspace
ring buffer. The collector normalizes them into typed events and
feeds disjoint per-process windows to the n-gram feature extractor.
When a window closes, the extractor emits a fixed-shape feature
vector. The scoring engine compares it against the cached baseline
for that process and raises an alert when the drift score exceeds
the per-process threshold. A SQLite store persists baselines and is
read at startup to warm the scorer's in-memory cache.

\subsection{Cross-language reference and production split}

\sediment{} and the n-gram extractor each have two implementations:
a Python reference in
\href{run:../chronosynd-py/}{\texttt{chronosynd-py/}} and a Rust
production port in
\href{run:../chronosynd-rs/}{\texttt{chronosynd-rs/}}. Python is
authoritative during research iteration. Rust is authoritative for
deployment. CI runs cross-implementation parity tests on every push,
covering both the baseline estimator (15 cases at three sample sizes
and four estimator configurations, 20 score inputs each) and the
syscall n-gram extractor (4 event-stream cases at varied window
sizes and vocabularies). Divergence between sides fails the build.

\subsection{Tamper-evident store}

Every state-changing operation (\texttt{put\_baseline},
\texttt{start\_maintenance\_window}, \texttt{end\_maintenance\_window})
commits in a single transaction together with an audit row. The
audit row's SHA-256 hash covers the operation, the payload, and the
previous row's hash. The chain is anchored at the all-zero hash, so
any tampering breaks the chain at the tampered row.
\texttt{verify-store} walks the chain and exits with a distinguished
code on tamper detection, so an operator's monitoring system can
tell tampering apart from a CLI crash.

\subsection{Operator surface}

Operators interact with \chronosynd{} through a single CLI binary.
The subcommands cover baseline inspection, tagging maintenance
windows when planned changes will shift the baseline, walking the
audit log to detect tampering, fitting baselines from CSV
observations or recorded JSONL event traces, and running the
collector daemon end-to-end against a configured store.
